VoiceLink
Helping deaf and hearing people communicate seamlessly in real time

Inspiration
Communication between deaf and hearing people can be slow and frustrating. We wanted to create a simple, accessible tool that bridges this gap in real time. Our inspiration came from seeing the difficulty some people face when trying to communicate quickly without interpreters or assistive devices.

What it does
VoiceLink allows a deaf user to communicate with a hearing person in three ways: -Typing: Users can type messages, which the app also speaks aloud. Hearing users then respond verbally, and VoiceLink converts their speech into text on the screen. -Hand gesture recognition: The user performs simple hand gestures in front of the laptop camera. The app recognizes these gestures and converts them to spoken English. -Quick Responses: The user can use the quick response feature to respond to a question in an instant without having to type. This creates a two-way, real-time conversation that is fast, intuitive, and inclusive.

How we built it
-Frontend web app built with React and plain JavaScript. -Camera input: Accessed via getUserMedia to detect predefined ASL gestures (Hello, How are you?, What’s your name?, Thank you, Goodbye). For the prototype, gesture detection uses simplified placeholder logic. -Text-to-Speech: Typed messages and recognized gestures are spoken using the browser’s Speech Synthesis API. -Speech-to-Text: Microphone input is captured via the Web Speech API and displayed as live captions for the deaf user. -UI: Two clear sections: “Sign or Type to Speak” and “Speak to Read” with high-contrast text and large buttons for accessibility.

Challenges we ran into
-Gesture detection: Accurately recognizing actual ASL in real time is complex. We simplified the detection for this hackathon using simple rules and placeholders in. -Browser compatibility: Some speech APIs behave differently across browsers; we focused on Chrome for reliable testing. -Timing and latency: Ensuring that speech and text updates felt instantaneous required careful handling of async APIs

Accomplishments that we're proud of
-Built a functioning prototype that supports two-way communication. -Implemented both camera-based hand gestures, typed messages, and quick responses in a single interface. -Created a clean, accessible UI suitable for real users with auditory accessibility needs. -Demonstrated social impact and technical feasibility.

What we learned
-Real-time accessibility tools require careful UX design to make them intuitive. -Simplifying gestures for prototypes can still demonstrate a strong impact without full ML integration. -Browser APIs like Speech Synthesis and Web Speech are powerful tools for inclusive apps.

What's next for VoiceLink
-Expand ASL coverage: Add more gestures for richer conversations and begin a transfer to true American Sign Language. -Integrate real ML models: Use hand-tracking and computer vision for accurate gesture recognition. -Mobile optimization: Make the app work on tablets and phones. -Customization: Allow users to add their own phrases and shortcuts for faster communication.

Built With
